% Generated draft: requires mathematical and source review.
\documentclass[11pt]{article}
\usepackage[margin=25mm]{geometry}
\usepackage{fontspec}
\setmainfont{Latin Modern Roman}
\usepackage{amsmath,amssymb,mathtools,unicode-math}
\setmathfont{Latin Modern Math}
\usepackage{enumitem,needspace,xurl,hyperref}
\setlength{\parindent}{0pt}
\setlength{\parskip}{5pt}
\setlength{\emergencystretch}{3em}
\newcommand{\N}{\mathbb{N}}
\newcommand{\Z}{\mathbb{Z}}
\newcommand{\Q}{\mathbb{Q}}
\newcommand{\R}{\mathbb{R}}
\newcommand{\C}{\mathbb{C}}
\newcommand{\F}{\mathbb{F}}
\newcommand{\A}{\mathbb{A}}
\newcommand{\PP}{\mathbb{P}}
\newcommand{\E}{\mathbb{E}}
\newcommand{\eps}{\varepsilon}
\newcommand{\dd}{\,\mathrm{d}}
\begin{document}
\section*{UnsolvedMath ID: 10400164}
\textbf{Source problem code:} catalogue entry
\subsection{Definitions and mathematical statement}
\emph{Problem (T. Kerler, nonsemisimple versus semisimple TQFTs and the double conjecture).} Retain the category-dependent theories denoted by \(V_{\mathcal C}\), \(W_{\mathcal B}\), and \(V_{D(\mathcal B)}\) in the source. Here \(D(\mathcal B)\) refers to the Drinfel'd double construction mentioned in the accompanying discussion. The supplied excerpt does not define the categories, the construction of these theories, or their admissibility hypotheses.
\begin{enumerate}
\item Clarify the difference in content between the two theories whose names both flatten to \(V_{\mathcal C}\) in the extracted record. Ask whether a homological TQFT \(H\) makes one essentially equivalent to the tensor product of the other with \(H\). The extracted comparison literally reads \(V_{\mathcal C}\) versus \(V_{\mathcal C}\otimes H\); missing typographical distinctions must be recovered before this can be an unambiguous equivalence assertion. In the usual meaning of homological, the mapping-class-group actions of \(H\) factor through their actions on first homology; the exact convention and meaning of ``essentially equivalent'' are not supplied.
\item Construct \(W_{\mathcal B}\) extending the Turaev--Viro TQFT construction to nonsemisimple \(\mathcal B\), analogously to the generalization of the Reshetikhin--Turaev construction in [216]. For quantum-group examples and closed \(3\)-manifolds, the construction should reproduce a version of the Kuperberg invariant. The version and normalization are not specified.
\item Determine the relationship between \(W_{\mathcal B}\) and \(V_{D(\mathcal B)}\); in particular, ask whether they are isomorphic as TQFTs in an appropriate sense. Neither this sense nor the full hypotheses on \(\mathcal B\) are provided.
\end{enumerate}
The source records genus-one evidence for the \(U_q(\mathfrak{sl}_2)\) case suggesting that the \(H\) in (1) is the Frohman--Nicas \(U(1)\) theory. It anticipates technical requirements in (2), as suggested by [236]: ``minimal'' cell decompositions of cobordisms rather than general triangulations, and combings rather than framings. It identifies (3) with Question 5 of [209], motivated by work of and discussions with D. Kazhdan and S. Gelfand in 1994, while allowing that it may have been posed earlier. For categories arising from subfactors and for closed manifolds, it reports results answering the conjecture in [204]. It presents more general results as a route to understanding the two TQFT constructions and a topological picture of the Drinfel'd double; no unrestricted equivalence is asserted as established.
\subsection{Short English statement}
Kerler asks how semisimple and nonsemisimple TQFTs differ, whether a homological factor relates them, how to extend Turaev--Viro theory to nonsemisimple categories, and whether that extension agrees with the theory of the Drinfel'd double. Some notation distinguishing the two theories in the first question is missing from the extraction.
\subsection{Sources}
\begin{itemize}
\item Local UnsolvedMath registry, permanent ID 10400164, source code catalogue entry. Attribution T. Kerler; source heading ``Non-semisimple vs. semisimple TQFT's, the double conjecture.'' The stored title is truncated after ``(T.''.
\item Ohtsuki, \emph{Problems on invariants of knots and 3-manifolds} (2002), Problem 8.18, PDF page 132. Supplied URL: \url{https://msp.org/gtm/2002/04/gtm-2002-04-024s.pdf}; recorded access 2026-07-29; extraction pdf-font-structure. The collection distinguishes 231 numbered bold problem, conjecture, and question headings.
\item Internal references [216] for the nonsemisimple Reshetikhin--Turaev generalization, [236] for anticipated technical issues, [209], Question 5, for the double conjecture and its discussion, and [204] for results for subfactor categories and closed manifolds. Full bibliographic entries are absent. D. Kazhdan and S. Gelfand are credited in the 1994 motivation.
\item The supplied registry labels the problem open and OPEN-TRIAGE. Current primary-literature status remains NEEDS\_REVIEW, as do redistribution terms for the publicly accessible author/publisher copy. The source's specialized results and genus-one evidence are retained with their stated scope.
\end{itemize}
\end{document}
